\chapter[Conditional Mechanics of the Graviton Record]{The Graviton Record Writer\\[0.2cm] \large Conditional Mechanics of a Persistent State-0 Strain Record}
\label{app:state0_record_writer}

\noindent\textbf{Abstract:}
Tension Field Theory identifies the graviton with the longitudinal strain wake accompanying an active photon and assigns the Hard Floor a persistent record of the interaction. This appendix examines a conditional reference-registry representation of that record. Integer site labels generate exactly compatible scalar bond differences, while a proposed event-dependent update connects the complete Tension-Column interface history to a change of registry. Two restrictions follow: a common temporal phase winding produces no internal longitudinal registry difference, and a continuous integer-valued selector is locally constant on connected families of admissible histories. These results constrain the candidate; they do not determine its physical transition law. The dimensional strain realization, reciprocal interaction, repeated-write response, and distinguishing readout remain to be supplied by the same parent mechanics. Definition~V persistence is preserved without introducing a Hard-Floor oscillator. The construction is a proposed, falsifiable representation, not a completed microscopic derivation.

\section{The Wake and the Written Configuration}
\label{sec:AW_wake_and_record}

\noindent\textbf{The photon moves. The graviton marks its passage.}
The coherent transverse excitation of the FCC Active Vacuum carries the photon action, while its Tension Column establishes the longitudinal strain wake identified with the graviton in \hyperref[app:graviton]{Appendix~\ref*{app:graviton}}. The induced wake follows the current loading; the written configuration retains a particular structural record after that loading has passed.

\hyperref[def:definition_five]{Definition~V} supplies the persistence requirement. It does not specify the event-dependent interaction that selects the retained pattern. This appendix separates that constitutional premise from one proposed representation and the mathematical restrictions that follow from it.

State~0 remains the foundational SC substrate, State~I the active FCC--condensate configuration, and State~II the compressed saturated phase. A State-II saturation transition is not automatically the writing process associated with an ordinary State-I photon. Nor is the active node--condensate coordinate in \hyperref[app:gravitational_endpoint_transfer]{Appendix~\ref*{app:gravitational_endpoint_transfer}} independently identified with State-0 memory merely because it is described as latent.

\noindent\textbf{Persistence is constitutional; writing requires a constitutive law.}
The task is to identify which longitudinal relationships change, how the existing record affects that change, and how the result couples to an active probe. No claim of a unique or exhaustive microscopic realization is made here.

\section{A Conditional Discrete Reference-Registry Representation}
\label{sec:AW_registry_representation}

Consider a finite, connected, contractible cubical patch of the SC reference network, with fixed connectivity and no periodic identifications. Let $V$ be its site set. Assign each site $v$ an integer label $n_v$ and collect the labels into $\mathbf n\in\mathbb Z^{|V|}$. These are proposed reference-registry branches, not physical displacements or additional oscillators.

For an oriented bond $b=(v\to w)$, define the site-to-bond difference operator $B$ by $(B\mathbf n)_b=n_w-n_v$. Let $F$ sum oriented bond differences around each elementary square face. The candidate normalized record is
\begin{equation}
    \mathbf e_{\mathrm R}=B\mathbf n,
    \qquad
    FB=0.
    \label{eq:AW_registry_gradient}
\end{equation}
Each site contribution cancels around a face. Consequently,
\begin{equation}
    F\mathbf e_{\mathrm R}=0.
    \label{eq:AW_compatible_record}
\end{equation}
This is exact scalar registry compatibility. Conversely, an integer bond field satisfying Equation~\eqref{eq:AW_compatible_record} can be integrated along lattice paths to obtain integer site labels: contractibility makes the integral independent of the chosen path, and connectedness leaves only a common additive offset. Additional prescribed boundary values must be checked separately. On a domain with holes or periodic identifications, face closure alone is insufficient; circulation around noncontractible loops must also vanish.\footnote{These are standard cellular cochain and path-integration properties. See A.~Hatcher, \emph{Algebraic Topology}, Cambridge University Press, 2002, Chapters~1--3.}

\noindent\textbf{A common relabeling is not a record.}
For any integer $c$,
\begin{equation}
    B(\mathbf n+c\mathbf 1)=B\mathbf n.
    \label{eq:AW_reference_gauge}
\end{equation}
The candidate record therefore depends on registry differences, not on a uniform offset. Setting one reference label to zero fixes a representative; it does not introduce another mechanical clamp. Comparisons at a physical patch boundary must include the bonds connecting it to the unchanged surrounding reference.

\medskip
\noindent\textbf{Registry compatibility is not yet elastic compatibility.}
The integer field $\mathbf e_{\mathrm R}$ is dimensionless. It is not a derived displacement field, strain tensor, stored energy, or mechanical equilibrium. A physical embedding must determine how a registry difference becomes longitudinal deformation while satisfying the SC geometry, constraints, and boundary tractions. Neither a displacement of $\ell_{\mathrm S}$ per integer step nor any other strain quantum is assigned by Equation~\eqref{eq:AW_registry_gradient}.

Because $B\mathbf n$ is an exact scalar difference, it has zero circulation around every closed lattice path in this patch. Integer labels alone therefore do not establish a dislocation, nonzero spatial vortex charge, or topological protection. The spatial phase winding discussed in \hyperref[app:physical_wavefunction]{Appendix~\ref*{app:physical_wavefunction}} is a separate construction. The registry candidate represents the constitutional record only if its branches correspond to physically distinct retained configurations.

\section{The Proposed Write and Update Law}
\label{sec:AW_write_update}

\noindent\textbf{The interface supplies the history; the transition selects the pattern.}
Let $\mathcal H_{\partial}$ denote the complete objective interface history of an event: the active photon mode, relative phase, Tension-Column loading, permitted contact deformations and rotations, and participating condensate and boundary response. The pre-existing registry is $\mathbf n^{-}$. In a physical completion, this history and the reciprocal active response must be solved together, not prescribed independently of the changing record.

The conditional update is
\begin{equation}
    \mathbf n^{+}
    =\mathbf n^{-}
    +\boldsymbol{\nu}[\mathcal H_{\partial},\mathbf n^{-}],
    \qquad
    \boldsymbol{\nu}\in\mathbb Z^{|V|}.
    \label{eq:AW_candidate_update}
\end{equation}
The selector $\boldsymbol{\nu}$ specifies a branch change, but Equation~\eqref{eq:AW_candidate_update} does not determine that selector. The resulting registry increment satisfies
\begin{equation}
    \Delta\mathbf e_{\mathrm R}=B\boldsymbol{\nu},
    \qquad
    F\Delta\mathbf e_{\mathrm R}=0.
    \label{eq:AW_update_compatibility}
\end{equation}
A nontrivial record change requires $B\boldsymbol{\nu}\neq0$, not merely $\boldsymbol{\nu}\neq0$. For a history expressed relative to the same physical reference, a consistent selector must also respect the uniform relabeling:
\begin{equation}
    B\boldsymbol{\nu}[\mathcal H_{\partial},\mathbf n+c\mathbf 1]
    =B\boldsymbol{\nu}[\mathcal H_{\partial},\mathbf n].
    \label{eq:AW_selector_reference_invariance}
\end{equation}

A localized site update changes its incident bond differences while preserving face closure. An independently assigned increment on one bond belonging to a square face generally does not. Thus even a localized imprint must be represented by a compatible collection of changes.

The objective FCC contact law supplies active traction, not an independently specified State-0 transition. In particular, the reference Hessian does not determine every nonlinear phase--deformation coupling, as shown in \hyperref[sec:D_phase_deformation_diagnostic]{Section~\ref*{sec:D_phase_deformation_diagnostic}}. Compatibility and tangent contact coefficients are constraints on a candidate writer, not substitutes for its selection law.

\medskip
\noindent\textbf{The next write encounters the existing record.}
Successive events require
\begin{equation}
    \mathbf n_{j+1}
    =\mathbf n_j
    +\boldsymbol{\nu}[\mathcal H_{\partial,j},\mathbf n_j].
    \label{eq:AW_repeated_writes}
\end{equation}
State-independent addition is only a special case. Whether writes commute, preserve their order, or alter the response of an already written address is determined by the selector and the coupled interface dynamics. Calibration on the unrecorded reference alone does not establish that dependence.

\section{Two Tests of the Registry Candidate}
\label{sec:AW_photon_tests}

\subsection{Common Photon Phase Is Not a Longitudinal Record}
\label{sec:AW_common_phase_test}

For a coherent traveling harmonic, the active phase is
\begin{equation}
    \varphi_v(t)=\mathbf k\cdot\mathbf r_v-\omega t.
    \label{eq:AW_coherent_phase}
\end{equation}
Test the provisional rule that each nonvanishing interface phasor $\psi_v=A_v e^{i\varphi_v}$ writes its temporal winding:
\begin{equation}
    w_v=\frac{1}{2\pi i}\oint\frac{d\psi_v}{\psi_v},
    \qquad
    \boldsymbol{\nu}=\mathbf w.
    \label{eq:AW_naive_phase_winding}
\end{equation}
The integer winding is defined here for a closed phasor path that never crosses zero. For fixed nonzero $A_v$ and one complete period $2\pi/\omega$, $d\psi_v/\psi_v=-i\omega\,dt$, so every participating site has $w_v=-1$. If the same increment is assigned to every site in the considered patch,
\begin{equation}
    \Delta\mathbf e_{\mathrm R}
    =B\mathbf w=B(-\mathbf 1)=0.
    \label{eq:AW_common_mode_cancellation}
\end{equation}
Spatial phase offsets do not change this result. The bond-relative phasor is also time independent for this harmonic, so its temporal winding vanishes.

\noindent\textbf{Action normalization does not select a registry branch.}
A selected active canonical mode can simultaneously satisfy
\begin{equation}
    \oint p\,dq=h,
    \label{eq:AW_complete_action}
\end{equation}
where $p$ is the momentum conjugate to $q$, not the photon's translational momentum. For example, $q=A\cos\theta$ and $p=-h\sin\theta/(\pi A)$ give Equation~\eqref{eq:AW_complete_action} exactly. Thus nonzero completed action coexists with zero registry change under Equation~\eqref{eq:AW_naive_phase_winding}. The action belongs to the complete excitation; the multimode normalization of \hyperref[app:photon_hamiltonian]{Appendix~\ref*{app:photon_hamiltonian}} is not assigned independently to each substrate address.

If an update instead assigns a common increment to a visited set $S\subset V$ and zero elsewhere, $B\mathbf 1_S$ is supported only on bonds crossing its boundary. Equal updates at adjacent visited sites cancel internally. This is an algebraic test of a patch-update rule, not a definition of phase winding at sites where the phasor vanishes.

\medskip
\noindent\textbf{A finite packet needs its own interface history.}
An envelope that vanishes before and after passage does not supply a nonvanishing closed phasor loop at each address. An arbitrary closing path or branch cut cannot supply the missing physical write event. Equation~\eqref{eq:AW_common_mode_cancellation} excludes the common-winding rule in its stated harmonic regime; it does not exclude an interface-generated, spatially selective record for a finite photon packet.

\subsection{Integer Selection and Physical Transition Boundaries}
\label{sec:AW_discrete_selection_test}

Fix the old registry and the event boundary conditions. If a family of admissible interface histories is connected and $\boldsymbol{\nu}$ is continuous on that family, its integer-valued image is a single point: the continuous image of a connected set is connected, whereas the integer lattice is discrete. Consequently, a continuously varying integer selector is constant within each such family. This does not force that constant to be zero unless a no-write history belongs to the same family.

\noindent\textbf{The discrete outcome can change while the motion remains continuous.}
Different outcomes require distinct history sectors or a boundary across which the integer outcome map is not continuous. That boundary need not be a singularity of the microscopic trajectory. Smooth dynamics can lead to different final branches when initial data cross a basin boundary. A physical writer must derive the branches and their selection boundary rather than assign a rounding threshold to an otherwise arbitrary continuous loading variable.

For the particular winding diagnostic, the integer is preserved under continuous deformations through closed, nonvanishing phasor loops. Changing it requires leaving that class, for example through a zero or loss of the endpoint closure. A single loop with nonzero winding need not itself pass through zero. Neither statement establishes a State-0 transition criterion.

These restrictions apply to the proposed integer-registry representation. They neither prove that every TFT record must be integer-valued nor exclude a persistent continuous reference configuration with a different constitutive structure.

\section{Static Retention, Work, Recoil, and Readout}
\label{sec:AW_physical_closure}

\noindent\textbf{Retention.}
In a fixed reference convention, the candidate expresses the post-event requirement as
\begin{equation}
    \mathbf n(t)=\mathbf n^{+}
    \quad\text{between further writing interactions}.
    \label{eq:AW_static_retention}
\end{equation}
Definition~V requires the physical longitudinal configuration to remain. Equation~\eqref{eq:AW_static_retention} represents that requirement; it does not derive the settling process or prove retention against an arbitrary later disturbance. No Hard-Floor kinetic term, oscillation frequency, or damping mechanism is introduced.

\medskip
\noindent\textbf{Writing work must be counted once.}
For a specified event and accounting region, the required energy balance is
\begin{equation}
    \Delta E_\gamma
    +\Delta E_{\mathrm{AV}\setminus\gamma}
    +\Delta E_{\mathrm{int}}
    +\Delta U_{\mathrm R}
    =W_{\mathrm{ext}}+\mathcal E_{\partial}.
    \label{eq:AW_energy_ledger}
\end{equation}
Here $E_{\mathrm{AV}\setminus\gamma}$ contains the other active degrees of freedom, including their recoil energy; $E_{\mathrm{int}}$ contains interaction energy not already included elsewhere; and $U_{\mathrm R}$ is the energy assigned to a physically realized retained configuration. The signed quantity $\mathcal E_{\partial}$ is net energy transported into the region, excluding the external work already counted in $W_{\mathrm{ext}}$. For an isolated system containing all participating fields, both right-hand terms vanish. For a fixed region crossed by the photon, its entering and leaving energy must be included in the transport balance. If all other energies return to their initial values in an isolated event, the balance forces $\Delta U_{\mathrm R}=0$; a positive retained-energy increment requires a compensating change elsewhere.

No function $U_{\mathrm R}(\mathbf n)$ is postulated here. A valid assignment must respect equivalent registry representatives, and its increment would be $U_{\mathrm R}(\mathbf n^{+})-U_{\mathrm R}(\mathbf n^{-})$, which can depend on the old record. The reduced $E_{\mathrm g}$ in \hyperref[app:graviton]{Appendix~\ref*{app:graviton}} is not automatically this writing work at every handoff. Likewise, repeating the address count does not multiply the available photon action or energy.

Energy conservation alone does not select the reciprocal force: forces differing by a component orthogonal to the instantaneous velocity have the same power. Momentum and angular-momentum balance, boundary impulses, and active recoil must therefore follow from the same interaction. The ideal Hard Floor is not an untracked finite-inertia recoil reservoir. An ideal zero-work constraint does not rule out energy-degenerate configurations, but neither does it supply their writing or readout law.

\medskip
\noindent\textbf{A final configuration is not automatically a complete archive.}
State-independent opposite increments cancel algebraically, and commuting additions lose event order. Neither property proves that physical erasure is permitted or that every distinguishable history is preserved. To represent an archive of successive traces, the admissible updates, distinct addresses, or a physically justified finer state must retain the required distinctions. The simple site labels do not establish that capacity.

\medskip
\noindent\textbf{The probe must couple to the retained configuration.}
Different $B\mathbf n$ are distinct candidate records, not automatically distinguishable physical states. An admissible probe must have a branch-dependent response determined by the actual strain embedding and coupling. A response depending only on a terminal active phasor cannot distinguish winding histories with identical final phasors. Reading must also include its own work, recoil, and possible alteration of the record. These requirements preserve the separate active measurement and coherent-accessibility constructions in \hyperref[app:deterministic_measurement_resolution]{Appendix~\ref*{app:deterministic_measurement_resolution}}, \hyperref[app:entanglement]{Appendix~\ref*{app:entanglement}}, and \hyperref[app:quantum_entanglement]{Appendix~\ref*{app:quantum_entanglement}}.

\section{Constitutive Completion and Falsification}
\label{sec:AW_completion}

\noindent\textbf{One interaction must connect the complete process.}
Physical completion requires a retained State-0 configuration with a defined strain embedding, a spatially selective update derived from objective interface dynamics, and an admissible probe response. The same law must operate on both an unrecorded reference and an already written substrate, with consistent action, work, impulse, and settling. A hypothetical update is not established merely because it satisfies Equation~\eqref{eq:AW_update_compatibility}.

The decisive propagation test is an ordinary coherent photon with its actual finite footprint. A candidate relying on an interface transition must show how that transition is generated under ordinary photon loading; regularity of the incoming carrier alone neither supplies nor forbids such a transition. A uniform temporal clock update must also confront Equation~\eqref{eq:AW_common_mode_cancellation} and the patch-boundary result. These tests must be met without fitting the selector to $E_{\mathrm g}$, a desired bit count, or an observational target.

The reference FCC tangent law, coherent condensate constraint, and ideal Hard-Floor compatibility relations do not independently specify the retained-state update. Their limits are maintained, rather than replaced by a new stiffness, barrier, or receiver. In particular, the singular active-coordinate reduction is not analytically continued into a latent evolution law; the distinction in \hyperref[sec:DMR_structural_limits]{Section~\ref*{sec:DMR_structural_limits}} remains in force.

\section{Conclusion: From Strain Wake to Persistent Information}
\label{sec:AW_conclusion}

The discrete registry candidate provides an exact scalar compatibility construction and a testable form for event-dependent updates. It also explains why a common photon phase turn is insufficient and why a physical branch-selection boundary must be distinguished from a mathematical label change. The selection, dimensional strain, reciprocal interaction, and readout are not derived by these results.

\noindent\textbf{The graviton is the strain wake. Its written configuration is the record. The interface mechanics must connect the two.}

The graviton identity and persistence belong to the adopted TFT architecture. The registry realization remains a conditional constitutive proposal, with its mathematical content separated from the physical mechanism still to be completed.
